% ---------------------------------------------------------------- 7 conclusion
\section{Conclusion}
\label{sec:conclusion}

This article presented doubly nested semantic communication, a token-level transceiver with a large-language-model prior that is nested in rate and in time and that charges every part of its operation in channel symbols. The rate-nested codebook with per-shell block power control makes every prefix a valid transmission and token-level incremental redundancy exact, without measurable loss at the prefix length that was examined. A Bayesian beam-search receiver requests further symbols only for uncertain tokens, and three receiver memories adapt its prior at three timescales, with every synchronized bit recorded and converted into channel symbols.

On a hierarchical drift stream with six seeds and a decoder selected on validation text, DN-SC reached 90\,\% token accuracy with 7.08 real channel symbols per token, against 9.23 for a frozen fixed-rate receiver with the same decoder. This is a reduction of 23.3\,\% with the return-link signaling charged, observed on every seed, and of 27.9\,\% (95\,\% CI 21.6--31.7\,\%) under free-running decoding of the audited messages. Most of the gain of the memories stems from the conversation context, which costs nothing on the air. The shipped memories lowered the prior by a further 0.37 bits per token but recovered their synchronization cost, relative to the context memory with refinement, only when one update served more than about 6100 messages. For a single link, the context memory with selective refinement is therefore the economical design, whereas shipped adaptation becomes worthwhile when one update serves a cell or a multicast group.

Future work will repeat the study with larger language models and on streams with naturally occurring drift, replace the capacity charge of signaling and synchronization by implemented codes on an error-prone return link with the delay of the rounds modeled, share one shipped update among the users of a cell, which is the regime in which the crossover analysis predicts that adaptation pays, and derive the shipping period of each memory from its size and drift rate, so that the schedule of the memories follows from the cost model rather than from a design choice.
